\subsection{Introducción Teórica}
\label{subsec:introduccion-teorica}

% ####################################################################################################################
% ####################################################################################################################

\subsubsection{Mediciones e Incertezas}

Toda medición experimental contiene una incerteza inherente.
Esta incerteza puede ser de naturaleza \textbf{sistemática} (siempre del mismo signo, debida a calibración del instrumento, método inadecuado, etc.) o \textbf{accidental/casual} (de signo variable, debida a errores de apreciación del observador, fluctuaciones del objeto, etc).

Para un conjunto de mediciones directas, se establece el siguiente criterio \cite{catedra2007mediciones} para determinar el valor representativo y la incerteza absoluta:
\begin{itemize}
    \item Valor medio:
    \begin{equation}
        X_0 = \frac{1}{2}(X_{M} + X_{m})
        \label{eq:valor-medio}
    \end{equation}
    \item Incerteza absoluta:
    \begin{equation}
        \Delta X = \frac{1}{2}(X_{M} - X_{m})
        \label{eq:incerteza-absoluta}
    \end{equation}
\end{itemize}

Siendo $\varepsilon_I$ la incerteza del instrumento y $V_F$ las variaciones fluctuantes del propio instrumento a medir.
Cuando $V_F < \varepsilon_I$, el instrumento no tiene la sensibilidad suficiente para detectar las irregularidades del objeto, por lo que se usa $\varepsilon_I$ directamente como incerteza absoluta.
Cuando $V_F \geq \varepsilon_I$, el instrumento es suficientemente preciso para capturar las irregularidades físicas, y se usa la fórmula antes mencionada.

La expresión final de una medición directa es: $X = X_0 \pm \Delta X$.

% ####################################################################################################################
% ####################################################################################################################

\subsubsection{Densidad de un Cilindro}

Se define la densidad $\rho$ de un cuerpo de masa $M$ y volumen $V$ como:
\begin{equation}
    \rho = \frac{M}{V}
    \label{eq:densidad}
\end{equation}

Usando $\varepsilon_x = \frac{\Delta X}{X_0}$ como definición de incerteza o error relativo, y aplicando propagación de errores mediante derivadas parciales \cite{baird1991experimentacion}.
Se calcula la incerteza relativa de la densidad como:
\begin{equation}
    \varepsilon_\rho = \varepsilon_M + \varepsilon_V
    \label{eq:incerteza-relativa-densidad}
\end{equation}

La deducción completa se desarrolla en el Apéndice~\ref{subsubsec:propagacion-errores-densidad}.

% ####################################################################################################################
% ####################################################################################################################

\subsubsection{Volumen de un Cilindro}

Se define el volumen $V$ de un cilindro de diámetro $d$ y altura $h$ como:
\begin{equation}
    V = \frac{\pi}{4} d^2 h
    \label{eq:volumen-cilindro}
\end{equation}

Análogamente, se calcula la incerteza relativa del volumen como:
\begin{equation}
    \varepsilon_V = 2\,\varepsilon_d + \varepsilon_h + \varepsilon_\pi
    \label{eq:incerteza-relativa-volumen-cilindro}
\end{equation}

La deducción completa se desarrolla en el Apéndice~\ref{subsubsec:propagacion-errores-volumen-cilindro}.

El error relativo introducido por $\pi$ ($\varepsilon_\pi = \frac{\Delta \pi}{\pi_0}$) proviene de recortar cifras de un número irracional.
Para minimizar $\varepsilon_\pi$ se elige la aproximación de $\pi$ que cumpla el criterio $10\,\varepsilon_\pi < 2\,\varepsilon_d + \varepsilon_h$.

Con los datos para el calibre, se adoptó $\pi = 3{,}142$.

Para la regla se adoptó $\pi = 3{,}14$, suficiente dada la menor precisión de ese instrumento.

El desarrollo del criterio de elección se puede encontrar en el Apéndice~\ref{subsubsec:criterio-eleccion-pi}.

% ####################################################################################################################
% ####################################################################################################################

\subsubsection{Péndulo Simple Ideal}

El péndulo simple ideal consiste en una masa puntual suspendida de un hilo inextensible y sin masa, que oscila sin fricción.
Para el régimen de \textbf{pequeñas oscilaciones} (donde se cumple la aproximación $\sen \theta \approx \theta$), el sistema se comporta como:

\begin{equation}
    T = 2\pi \sqrt{\frac{L}{g}} \quad \Longrightarrow \quad T^2 = \frac{4\pi^2}{g}\,L
    \label{eq:periodo-pendulo-angulos-pequennos}
\end{equation}

Despejando $g$ de la ecuación y propagando errores obtenemos:

\begin{equation}
    \varepsilon_g = \frac{\Delta g}{g_0} = \frac{\Delta L}{L_0} + 2 \cdot \frac{\Delta T}{T_0}
    \label{eq:error-relativo-periodo-pendulo-angulos-pequennos}
\end{equation}

Dado que la incerteza relativa introducida por $\Delta T$ es 20 veces mayor al necesario para cumplir con la incerteza requerida y siendo $\Delta T = \frac{\Delta t}{n}$, tomamos $n = 20$ oscilaciones para reducir el error.

El modelo de péndulo ideal supone que la masa en suspensión es puntual.
Esto no es cierto, ya que en la experiencia empleamos esferas de aproximadamente 1,7 $cm$ de radio.
Una corrección en la fórmula para el período que tenga en cuenta una masa esférica de radio r es:

\begin{equation}
    T = 2\pi \sqrt{\frac{L + \frac{2r^2}{5L}}{g}}
    \label{eq:periodo-pendulo-angulos-pequennos-corregido}
\end{equation}

Para el régimen de \textbf{grandes oscilaciones} (donde no se cumple la aproximación $\sen \theta \approx \theta$) una aproximación más general se logra tomando los primeros cuatro términos de la serie de Taylor:

\begin{equation}
    T \approx 2\pi \sqrt{\frac{L}{g}} \left[ 1 + \frac{1}{4} \sen^2 \left( \frac{\theta_0}{2} \right) + \frac{9}{64} \sen^4 \left( \frac{\theta_0}{2} \right) + \frac{225}{2304} \sen^6 \left( \frac{\theta_0}{2} \right) \right]
    \label{eq:periodo_real}
\end{equation}
